\documentclass[11pt]{article}

\usepackage{amsmath,amssymb}
\usepackage{xurl}
\usepackage[hidelinks]{hyperref}
\setlength{\emergencystretch}{2em}

% Shared commands for notes in docs/paper-gaps/.
%
% Two halves.  The link commands below are project identity: OWNER/REPO is the
% placeholder that scripts/bootstrap_project.py replaces with this project's
% GitHub slug and Pages site.  Everything after them is NOTATION, and notation
% belongs to the source paper: the definitions here are an example set, and a
% project replaces them with the macros of its own paper's style file.  The one
% rule that never changes: a command defined here must mean exactly what the
% source paper's macro means, or a note silently says something else.

\newcommand{\Lean}{\textsc{Lean}}
\newcommand{\leanid}[1]{\textsf{#1}}
\newcommand{\ghissue}[1]{\href{https://github.com/OWNER/REPO/issues/#1}{GitHub issue \##1}}
\newcommand{\ghpr}[1]{\href{https://github.com/OWNER/REPO/pull/#1}{PR \##1}}
% Local issue tracker (issues/NNNN-*.md in this repository); no remote URL.
\newcommand{\localissue}[1]{issue \##1 (\path{issues/})}
\newcommand{\doclink}[1]{\href{https://OWNER.github.io/REPO/docs/#1.html}{\path{#1}}}

% --- Notation: replace with the source paper's own macros --------------------
% \F, \N, \C, \R, \tr, \ot, \eps, \E, \bx, \bu, \bv, \by

\newcommand{\F}{\mathbb{F}}
\newcommand{\N}{\mathbb{N}}
\newcommand{\C}{\mathbb{C}}
\newcommand{\R}{\mathbb{R}}
\newcommand{\tr}{\mathrm{tr}}
\newcommand{\eps}{\epsilon}
\newcommand{\ot}{\otimes}
\newcommand{\E}{\mathop{\bf E\/}}

% bold random variables
\newcommand{\bu}{\boldsymbol{u}}
\newcommand{\bv}{\boldsymbol{v}}
\newcommand{\bx}{{\boldsymbol{x}}}
\newcommand{\by}{\boldsymbol{y}}

% T1 so literal < and > in \texttt placeholders typeset as themselves
\usepackage[T1]{fontenc}

% bra-ket
\usepackage{braket}

% --- Example structured spaces ---
\newcommand{\polyfunc}[3]{\mathcal{P}(#1, #2, #3)}
\newcommand{\polymeas}[3]{\mathrm{PolyMeas}(#1, #2, #3)}
\newcommand{\polysub}[3]{\mathrm{PolySub}(#1, #2, #3)}

% Approximation relations (paper uses \approx_\eps and \simeq_\zeta)
\newcommand{\approxeps}{\approx_{\varepsilon}}
\newcommand{\simgeq}{\simeq_{\zeta}}


\newcommand{\ind}{\operatorname{ind}}
\renewcommand{\ghissue}[1]{%
  \href{https://github.com/Dengnifer/MIPStarRE-QPBT/issues/#1}{GitHub issue \##1}}

\title{Corrections to the Pauli-Test Winning Implications}
\author{}
\date{2026-09-23}

\begin{document}
\maketitle

\section{At a glance}

\textbf{Difficulty.} The phase and index corrections are algebraic; the
conditioning argument uses an exact finite probability. The substantive
difference is the evaluation and acceptance convention on zero-direction
lines. The existing winning implications are proved for that convention;
their identification with the printed source convention is not established.

\textbf{Estimated weight.} This revision documents existing results and their
scope; it adds no formal proof. \textbf{Mathlib/project split.} Finite sums,
post-processing of measurements, and state-dependent norm estimates provide
the general infrastructure. The question distribution, completed evaluation,
and downstream constructions are project-specific.
\textbf{Key mathematical inputs.} Exact commuting and anticommuting probabilities, rejection bounds,
and a value-only Magic Square anticommutator estimate on the original state.
The verdict is a \emph{documented deviation}, not a proof or a refutation of
the unrestricted printed line-evaluation claim.%
\footnote{The original implication results are tracked in \ghissue{112} and
\ghissue{113}; this documentary disposition is \ghissue{714}.}

\section{Key theorem forms}

Let $q=2^k$ with $k$ odd, let $m,d\geq1$ with $m\mid q$, and let a projective
strategy on a unit bipartite state $\ket\psi$ succeed with probability at least
$1-\eps$, where $\eps\geq0$, in the formal Pauli basis game. Here that game
uses the coefficient-answer and universal-evaluation convention described
below. The trace $\tr$ takes values in $\F_2$, and the interpolation vector
$\ind_m(u)$ has coordinates
$\prod_{i:b_i=0}(1-u_i)\prod_{i:b_i=1}u_i$ indexed by
$b\in\{0,1\}^m$. For a uniform tuple
$\omega=(u_X,u_Z,r_X,r_Z)\in(\F_q^m)^2\times\F_q^2$, put
\[
  \gamma(\omega)
  =\tr\big((r_X\ind_m(u_X))\cdot(r_Z\ind_m(u_Z))\big).
\]
The corrected probabilities are
\[
  \Pr[\gamma\ne0]=\frac12(1-q^{-1})^{m+1}\geq\frac1{16},
  \qquad \Pr[\gamma=0]\geq\frac12.
\]
Thus the subtest consistency defects are $O(\eps)$, including those
conditioned on $\gamma=0$ or $\gamma\ne0$. The induced Magic Square strategy
has conditional average value $1-O(\eps)$ on $\gamma\ne0$.

For a sampled line--point pair $(\ell,u)$, let $L_a^{\ell,u}$ denote the
completed line-evaluation effect and let $B_a^u$ denote the point effect,
extended by $B_\bot^u=0$. The low-degree implication has the form
\[
  \E_{(\ell,u)}\sum_{\substack{a,b\in\F_q\cup\{\bot\}\\a\ne b}}
    \bra\psi L_a^{\ell,u}\ot B_b^u\ket\psi\leq C\eps
\]
for a universal $C$. The squared-distance companion compares the same
completed families on $\ket\psi$ and also has error $O(\eps)$.
For the observables
$W^r(u)=\sum_{a\in\F_q}(-1)^{\tr(ar)}M^{(\mathrm{Point},W),u}_a$
of the complete typed point measurements, one has self-consistency with
squared-distance error $O(\eps)$ and
\[
  X^{r_X}(u_X)Z^{r_Z}(u_Z)\ot I
  \approx_{\sqrt\eps}
  (-1)^{\gamma(\omega)}Z^{r_Z}(u_Z)X^{r_X}(u_X)\ot I
\]
on average over all tuples, with both player orientations. Here
$A\approx_\eta B$ means $\E\|(A-B)\ket\psi\|^2=O(\eta)$.
None of these established estimates assumes $6md\leq q$.

\section{The source assertions and discrepancies}

The source states the seven measurement implications in
\path{lem:qld-win-implications} and the observable implications in
\path{len:qld-win-implications-obs} of the Pauli analysis
\cite[Appendix~A, Strategies]{Ji2020MIPStar}.%
\footnote{See
\path{references/qpbt-paper/14_analysis_of_the_pauli_basis_test.tex},
lines~197--265 and 267--362.}
The hypotheses are success in the Pauli basis game for admissible parameters;
projectivity has already been obtained by dilation. The conclusions concern
question self-consistency, low-degree and Pauli-basis consistency, the two
commutation checks, and Magic Square value and consistency. The source's
low-degree item sums only over $\F_q$, without a missing outcome.

First, the displayed conditioning event for the Magic Square average omits
$\ind_m$ and writes $\tr((r_Xu_X)\cdot(r_Zu_Z))$. This expression uses the
dot product in $\F_q^m$, whereas the game's phase uses the indicator vectors
in $\F_q^{2^m}$. These events need not agree: at the admissible tuple
$(q,m,d)=(2,1,1)$, take $u_X=u_Z=0$ and $r_X=r_Z=1$. The printed expression
is zero, but $\ind_1(0)=(1,0)$ and $\gamma=1$. This refutes the equality of
the conditioning events within the source domain; it does not refute all
seven winning implications.

Second, the proof assumes $6md\leq q$ and disposes of the complementary range
by noting that the final Pauli soundness error can then be made trivial.
Such a reduction can restrict the part of a final theorem's proof requiring
this lemma, but cannot prove the lemma as an unrestricted standalone result.
The replacement uses
\[
  \ind_m(u_X)\cdot\ind_m(u_Z)
  =\prod_{i=1}^m(1+u_{X,i}+u_{Z,i}).
\]
Each factor is independently nonzero with probability $1-q^{-1}$ as $u_Z$
varies. Independently $r_X\ne0$ with that probability, and then the trace is
balanced as $r_Z$ varies. This is the established exact probability calculation,
which gives the bounds above since $m\mid q$ implies $m\leq q$.%
\footnote{The probability correction is documented separately in
\path{docs/paper-gaps/qpbt_anticommuting-probability.tex}; the declarations
are \leanid{anticommProb\_eq}, \leanid{commProb\_ge\_half}, and
\leanid{anticommProb\_ge\_of\_m\_le\_q} in
\path{MIPStarRE/QPBT/Observables/Anticommuting.lean}.}
The degree-zero example $(q,m,d)=(2,1,0)$, where the printed probability
bound fails, is outside the source domain: the paper uses positive natural
numbers. It is not a counterexample to the lemma on admissible parameters.

Finally, a display in the observable proof writes
$Z^{r_X}(u_X)X^{r_Z}(u_Z)$ where the intended reversed product is
$Z^{r_Z}(u_Z)X^{r_X}(u_X)$. The latter agrees with the statement of the
observable lemma and the variable indices in the same display. The source label prefix
\path{len:qld-win-implications-obs} is likewise a typographical error for
\path{lem:qld-win-implications-obs}.

\section{Evaluation and the actual game convention}

The source explicitly permits a line $\ell(u_0,v)=\{u_0+tv:t\in\F_q\}$
with $v=0$. It specifies coefficient lists as answers, but its low-degree
decider checks $f(t)=a$ for ``the'' parameter with $u=u_0+tv$. At $v=0$
and $u=u_0$, every parameter satisfies this equality. A nonconstant function
$t\mapsto f(t)$ therefore has no value determined by the point alone.%
\footnote{See \path{references/qpbt-paper/08_classical_and_quantum_low_degree_tests.tex},
\path{def:line}, lines~107--121, and \path{fig:ld-decider}, lines~337--386.
The Pauli low-degree check invokes that decider at lines~1148--1159.}
Such questions have positive probability. Conditional on a diagonal question,
their probability is $m^{-1}\sum_{j=1}^m q^{-j}$; the equal line--point mixture
has half that mass. Choosing a parameter, discarding these questions, and
requiring equality at every parameter are different conventions.

The formal game chooses the last convention. For $u\in\ell$, an answer
evaluates to $a$ precisely when
\[
  f(t)=a\quad\hbox{for every }t\in\F_q\hbox{ such that }u=u_0+tv.
\]
For nonzero direction there is a unique such $t$. At zero direction this
condition says that the induced function on $\F_q$ is constant, not that its
coefficient representative is a constant polynomial. These distinctions matter
when the degree bound is at least $q$. On pairs outside the sampled support,
the win predicate is vacuous if no parameter exists, whereas partial evaluation
is undefined there; every sampled point lies on its line.%
\footnote{The exact definitions are \leanid{dlinePointCondition} and
\leanid{ldWinPredicate} in \path{MIPStarRE/QPBT/Test/LowDegreeGame.lean},
and \leanid{DegPoly}, \leanid{EvaluatesTo}, and \leanid{evalOpt} in
\path{MIPStarRE/QPBT/Observables/LineDefs.lean}.}

For a complete line measurement $\{L_f\}$, partition the answers that evaluate
by their values, and put the remaining effects into $L_\bot$:
\[
  L_a=\sum_{f\text{ evaluating to }a}L_f,\qquad
  L_\bot=\sum_{f\text{ without an evaluation}}L_f.
\]
Then $\sum_{a\in\F_q}L_a+L_\bot=I$. This is a complete measurement indexed
by $\F_q\cup\{\bot\}$, whereas
the field-valued classes alone can have total effect strictly below $I$.
For example, at $(q,m,d)=(2,1,1)$, place effect $I$ on the allowed answer
$f(T)=T$ at a zero-direction diagonal question. It has no evaluation at the
sampled point, so $L_\bot=I$ and every $L_a=0$. The universal check rejects
against either field answer, while selecting $t=0$ would accept answer zero.
This is a local answer/measurement example on admissible parameters, not a
high-success strategy refuting the source lemma.

The point measurement is extended by $B_\bot=0$. Completion adds exactly
$\bra\psi L_\bot\ot I\ket\psi$ to the field-only off-diagonal consistency
defect, and adds $\|(L_\bot\ot I)\ket\psi\|^2$ to the corresponding
squared-distance sum. Thus completion does not merely rename an alphabet.
Every mismatch of completed labels is rejected by the universal check, so its
Born mass is controlled by the game's rejection probability. The extra symbol
is an internal measurement label, not a new permitted answer to the verifier.

There is a separate completion in the formal representation of answers. A
strategy measures one common sum-type alphabet, and malformed answers are
rejected. The typed point and line families are obtained by sending malformed
answers to the zero value or zero coefficient list; axis lists are padded to
degree $md$. Their measurements are complete, and their malformed-answer mass
is bounded by rejection, not assumed zero. Only after this post-processing is
the line evaluation completed by $\bot$; the point's $\bot$ effect is zero.%
\footnote{See the \leanid{ProjectiveSetting} methods \leanid{pointMeas},
\leanid{lineMeas}, and \leanid{lineEvalMeas} in
\path{MIPStarRE/QPBT/Observables/Defs.lean}, and
\leanid{pointMeasOption} in
\path{MIPStarRE/QPBT/Observables/WinImplications/Setup.lean}.}

Under this convention, conditioning a nonnegative rejection cost on the
commuting or anticommuting event costs at most $2$ or $16$, respectively.
The finite type graph contributes another universal factor. This gives the
seven established implications and their distance companions. For observables,
the formal proof uses the value-only Magic Square anticommutator bound on the
original state, for both players, and transfers it using point--variable
consistency. It does not require the unrestricted extraction claim of the
printed Magic Square rigidity theorem. That theorem has its own deviation
note. The resulting observable error is $O(\sqrt\eps)$ for all $\eps\geq0$.

\section{Blueprint and formalization status}

The blueprint's winning-implication nodes state the corrected phase and products
and the completed line comparison. Its remarks on coefficient representatives
and zero directions specify the game convention. The statement and proof
marks there concern those explicitly stated objects; they do not certify
equality with an unspecified source acceptance rule on zero directions.
The Lean hypotheses retain admissibility, projectivity, and success, with
$\eps\geq0$ explicit. No intermediate estimate is supplied as an extra premise.
The comparison differs in its evaluation alphabet and game convention, not in
an additional hypothesis such as $6md\leq q$.%
\footnote{The blueprint labels are \path{lem:qld-win-implications},
\path{lem:qld-win-implications-obs}, and
\path{rem:qld-win-implications-typos}. The seven \path{win_*} declarations,
their distance companions, \leanid{pointObs\_self\_consistent}, and both
\leanid{pointObs\_twisted\_commutation} orientations are proved in
\path{MIPStarRE/QPBT/Observables/WinImplications.lean} and its imported
leaves. The value-only estimates are
\leanid{msVarObs\_anticommutator\_le} and
\leanid{msVarObsA\_anticommutator\_le} in the
\path{WinImplications/AnticommutingObs.lean} leaf.}

In the source, these implications are used to expand the Hilbert spaces,
combine measurements, and extract the Pauli registers. The implemented route
uses the winning estimates with the following explicit
conventions. Tensoring with ancillary Pauli operators cancels the signed phase
and gives approximately commuting expanded point observables. The completed
line estimate gives consistency of the expanded line measurement, with every
opposite register placement proved directly. Field-valued point measurements
are then combined and rounded on the existing expanded spaces, without the
source's unproved ancillary-padding identification. The combined and extended
line constructions retain completed evaluations; at extended dimension $2m+2$
they use the directly indexed line law. The established extended-line error is
$C m\,\operatorname{poly}(\eps,md/q)$, not the printed
$\operatorname{poly}(m^2\eps,md/q)$ claim.%
\footnote{See \path{MIPStarRE/QPBT/Observables/ExpandedCommutation.lean},
\path{MIPStarRE/QPBT/Observables/LineMeasurement.lean},
\leanid{exists\_combinedPointsWitness} in
\path{MIPStarRE/QPBT/Combining/Points.lean},
\leanid{exists\_combinedLinesWitness} in
\path{MIPStarRE/QPBT/Combining/Lines.lean}, and
\leanid{exists\_extendedLinesWitness\_established} in
\path{MIPStarRE/QPBT/Combining/Apply.lean}. The separate differences are
documented in the combined-points, dimension-divisibility, combined-lines,
and subline-claims notes in \path{docs/paper-gaps/}.}

That established construction supplies a global polynomial-pair measurement
and hence the extraction witness. Point--Pauli consistency also enters the
nonencoding-mass and evaluated-Pauli comparisons. The final isometry and
completed-to-prescribed-answer transfers give the public soundness conclusions.
None of this proves a transport to the source's extended seed-indexed law, a
field-only POVM identification, or the printed extended-line error bound.%
\footnote{See \leanid{exists\_globalPairWitness} in
\path{MIPStarRE/QPBT/Combining/Apply.lean},
\path{MIPStarRE/QPBT/Extraction/EvaluatedPauliConsistency.lean},
\leanid{exists\_extractionWitness} in
\path{MIPStarRE/QPBT/Extraction/SourceUnitary.lean}, and
\path{MIPStarRE/QPBT/Test/Soundness/RawOperatorTransfer.lean}.
The public theorems are \leanid{pauli\_soundness} and
\leanid{pauli\_soundness\_qubit}; their statements are unchanged.}

\section{Conclusion}

The algebraic and probability repairs are established on the admissible domain.
The completed winning implications and their observable consequences are also
proved for the actual formal game. The coefficient, evaluation, and
zero-direction acceptance conventions are substantive differences whose full
identification with the printed source remains unresolved. The verdict
\emph{documented deviation} closes this intermediate documentation item only:
it neither asserts the unsupported source claim nor adopts a new game, changes
a theorem statement, or settles the separate downstream source-law questions.

\bibliographystyle{alpha}
\bibliography{references}

\end{document}
